\documentclass[11pt,a4paper]{article}

% ── Packages ─────────────────────────────────────────────────────────────────
\usepackage[T1]{fontenc}
\usepackage[utf8]{inputenc}
\usepackage{mathptmx}
\usepackage{amsmath,amssymb,amsthm}
\usepackage{textcomp}
% mathptmx drops \hbar — restore it from the amsfonts glyph
\usepackage{amsfonts}
\AtBeginDocument{\let\hbar\relax
  \DeclareMathSymbol{\hbar}{\mathord}{AMSb}{"7E}}
\usepackage{geometry}
\usepackage{microtype}
\usepackage{hyperref}
\usepackage{booktabs}
\usepackage{array}
\usepackage{enumitem}
\usepackage{setspace}
\usepackage{titlesec}
\usepackage{abstract}
\usepackage{graphicx}
\usepackage{caption}
\usepackage{natbib}

\geometry{top=1.1in, bottom=1.1in, left=1.15in, right=1.15in}

% ── Section heading style ─────────────────────────────────────────────────────
\titleformat{\section}{\normalfont\bfseries\large}{\thesection}{1em}{}
\titleformat{\subsection}{\normalfont\bfseries\normalsize}{\thesubsection}{1em}{}
\titleformat{\subsubsection}{\normalfont\bfseries\itshape\normalsize}{\thesubsubsection}{1em}{}
\titlespacing*{\section}{0pt}{14pt}{4pt}
\titlespacing*{\subsection}{0pt}{10pt}{3pt}
\titlespacing*{\subsubsection}{0pt}{8pt}{2pt}

% ── Abstract style ────────────────────────────────────────────────────────────
\renewcommand{\abstractnamefont}{\normalfont\bfseries}
\renewcommand{\abstracttextfont}{\normalfont\small}
\setlength{\absleftindent}{0.5cm}
\setlength{\absrightindent}{0.5cm}

% ── Theorem environments ──────────────────────────────────────────────────────
\theoremstyle{plain}
\newtheorem{theorem}{Theorem}[section]
\newtheorem{lemma}[theorem]{Lemma}
\theoremstyle{definition}
\newtheorem{definition}[theorem]{Definition}
\theoremstyle{remark}
\newtheorem{remark}[theorem]{Remark}

% ── Hyperref ─────────────────────────────────────────────────────────────────
\hypersetup{
  colorlinks=true, linkcolor=black, citecolor=black, urlcolor=black,
  pdftitle={The Relational Universe: Interaction and the Transition of
            the Possible to the Actual},
  pdfauthor={Joshua F. Sandeman}
}

% ── Macros ───────────────────────────────────────────────────────────────────
\newcommand{\ket}[1]{|#1\rangle}
\newcommand{\bra}[1]{\langle #1|}
\newcommand{\braket}[2]{\langle #1 | #2 \rangle}
\newcommand{\expect}[1]{\langle #1 \rangle}
\newcommand{\rhoS}{\rho_{S}}
\newcommand{\inst}{\mathcal{J}}          % quantum instrument
\newcommand{\OA}{\hat{O}_{A}}
\newcommand{\OB}{\hat{O}_{B}}

% ─────────────────────────────────────────────────────────────────────────────
\begin{document}

% ── Title block ───────────────────────────────────────────────────────────────
\begin{center}
  {\LARGE\bfseries The Relational Universe: Interaction and the\\[4pt]
   Transition of the Possible to the Actual\par}
  \vspace{10pt}
  {\large Joshua F.\ Sandeman$^{1*}$\par}
  \vspace{4pt}
  {\normalsize $^{1}$Independent Researcher, Salem, OR, USA.\par}
  {\normalsize $^{*}$Corresponding author: sansuikyo@gmail.com\par}
  \vspace{8pt}
  {\small\textit{Preprint:}~\href{https://doi.org/10.5281/zenodo.19174380}{doi.org/10.5281/zenodo.19174380}
  (concept DOI: \href{https://doi.org/10.5281/zenodo.19174379}{doi.org/10.5281/zenodo.19174379})\par}
  \vspace{14pt}
\end{center}

% ── Abstract ─────────────────────────────────────────────────────────────────
\begin{abstract}
We argue that quantum randomness is a logical necessity arising from the nature of
intrinsically relational properties---properties lacking definite states outside
specific contexts of causal interaction. This motivates Relational Actualism
(RAQM), an ontological framework classifying physical properties as invariant,
classical, or intrinsically relational. RAQM grounds wave function actualization
in physical interaction rather than observer-relative facts. The actualization
criterion is physically grounded and mathematically precise: an interaction
constitutes an actualization event when it produces an irreversible increase in
the quantum relative entropy of the system with respect to the vacuum reference
state, $\Delta S(\rho\|\sigma_0) > 0$, generating a robust, redundant causal
record. The on-shell condition $p^{\mu}p_{\mu} = m^{2}c^{2}$ is the
implementation of this entropic criterion in weakly coupled perturbative QFT;
it is not the fundamental definition and does not constitute a claim of sharp
spacetime localization (see Section~\ref{sec:localization} for the explicit
response to Reeh-Schlieder and LSZ concerns).
This criterion is basis-independent, non-perturbative, and requires no new
fundamental constants. To reconcile a globally consistent wave function with
special relativity, the Tomonaga-Schwinger equation establishes a
foliation-independent global update mechanism. We derive a consistency condition
demonstrating that RAQM's actualization map preserves identical observable
statistics at spacelike-separated points regardless of hypersurface choice---a
derived result, not a postulate---strictly forbidding superluminal signaling via
algebraic microcausality. The no-signaling result applies rigorously to the
non-selective (outcome-averaged) actualization map; conditioned branches are
treated explicitly. Three specific, discriminating empirical predictions are
stated: a finite actualization propagation velocity bounded by the Lieb-Robinson
speed, an environment-dependent transition duration scaling as $\hbar/\Delta
E_{\mathrm{env}}$, and a sharp reversibility threshold at the actualization
criterion. RAQM is distinguished from Relational Quantum Mechanics (which
relativizes facts), dynamical collapse models (which ground collapse in
stochastic gravitational thresholds rather than causal interaction),
decoherence-only models (which leave the outcome problem open), and Everettian
Many-Worlds (which denies definite events). Classical behavior emerges dynamically
when continuous actualization events eliminate all unconstrained degrees of
freedom. The companion papers RAGC and RACI extend these results to gravity,
cosmology, and biological complexity; RAHC extends
the framework to the emergence of agentic intelligence. Recent results:
the Einstein field equations are the unique consistent macroscopic description
of the RA causal graph by Lovelock's theorem;
$c = l_P/t_P$ is derived as the maximum causal bandwidth rate and
$E = \gamma mc^2$ from discrete counting alone; the Causal Firewall
threshold is proved via the Erd\H{o}s-R\'{e}nyi percolation theorem;
and the Poisson-CSG generative measure makes the causal graph
dynamics local and environment-dependent for the first time in any
sequential-growth formulation. Key theorems have been formally verified
in Lean~4/Mathlib. A second Lean~4 file
(\texttt{RA\_AQFT\_Proofs\_v10.lean}) formalises the finite-dimensional AQFT
proof program and proves the frame-independence and Rindler stationarity
theorems; see the note on formal verification in the companion paper
RAGC.

\medskip\noindent
\textbf{Keywords:} Relational Quantum Mechanics \textperiodcentered{} Wave
Function Collapse \textperiodcentered{} Microcausality \textperiodcentered{}
Decoherence \textperiodcentered{} Determinism \textperiodcentered{} Phase Space
\textperiodcentered{} Objective Collapse \textperiodcentered{} Entanglement
Entropy \textperiodcentered{} Lean~4 Verification \textperiodcentered{}
Causal DAG \textperiodcentered{} Actualization Criterion \textperiodcentered{}
Kinematic Snap \textperiodcentered{} Poisson-CSG
\end{abstract}

\hrule
\vspace{8pt}

% ─────────────────────────────────────────────────────────────────────────────
\section{Introduction}

If the universe is, at least under certain circumstances, random, then it is
impossible, even in principle, to build a deterministic model of it. In the best
case, it can yield possible outcomes and their probabilities of occurring. What
actually happens must be added to the model. As far as we know, there are indeed
circumstances in which the universe is truly random. Furthermore, randomness in
the universe is not some vexing and curious add-on; it is, in fact, logically
necessary---a consequence, as we shall see, of the relational nature of the
universe.

The paper proceeds as follows. Section~\ref{sec:landscape} situates the view
philosophically, distinguishing it from other interpretations including Relational
Quantum Mechanics, dynamical collapse models, and decoherence-based approaches.
Section~\ref{sec:formalism} develops the mathematical formalism in terms of
density matrices, decoherence, and the reduced density matrix.
Section~\ref{sec:experiment} walks through the key experimental evidence.
Section~\ref{sec:discussion} discusses the broader implications for determinism,
causality, the reality of the wave function, and empirical falsifiability.
Section~\ref{sec:conclusion} concludes.

% ─────────────────────────────────────────────────────────────────────────────
\section{Rethinking Quantum Theory: Interpretive Landscape}
\label{sec:landscape}

Erwin Schr\"{o}dinger once said of Quantum Theory, ``I don't like it, and I'm
sorry I ever had anything to do with it.'' What Schr\"{o}dinger objected
to---which inspired his famous thought experiment involving a hapless cat---was
randomness. Einstein shared his concerns. Quantum Theory has the dubious
distinction of being one of the most powerful and accurate models of the world
ever developed, while also being the source of considerable controversy,
bewilderment, and confusion. The root of the problem is that, under certain
circumstances, Nature appears to be truly random, which necessitates a significant
shift in thinking about how it can be modeled and how it evolves.

\subsection{A Taxonomy of Properties}

Phenomena in nature have a set of possible states, only one of which can be
actualized (manifest, measurable) at any given point in spacetime. Properties can
be categorized as follows:

\textbf{Invariant properties} are those that can only ever have one possible
state---e.g., the charge of an electron. Their probability of measurement in the
one state is unity.

\textbf{Classical properties} are either fixed or vary deterministically as the
result of deterministic interactions. They are actualized as a state of being. The
trajectory of a planet orbiting in a gravitational field is a paradigm case.

\textbf{Intrinsically relational properties} cannot assume definite values outside
a specific context of interaction. Examples include the location or spin of newly
created particles whose spins are correlated. Crucially, after undergoing an
appropriate interaction, intrinsically relational properties can become
classical---they become actualized and behave the same as invariant properties, so
long as the context of interactions eliminates all other possibilities.

The Special and General Theory of Relativity introduced the world to the concept
of relational properties, demonstrating that there is no absolute coordinate
system for spacetime. Quantum mechanics goes further, showing that intrinsically
relational properties, prior to an interaction that restricts the state space, are
multivalued. In Quantum Theory this is modeled as a linear combination of
amplitudes which, via the Born Rule, correspond with very high accuracy to the
probability that a certain state will be actualized given a particular type of
interaction.

\subsection{Positioning: Relational Actualism vs.\ Relational Quantum Mechanics}

Carlo Rovelli's Relational Quantum Mechanics (RQM) shares with the present
framework the conviction that quantum states are not absolute but are defined
relative to an observer or reference system \citep{Rovelli1996}. In RQM, there
are no observer-independent facts about the values of physical quantities; every
fact is a fact relative to some system.

The present view, Relational Actualism, agrees that quantum properties are
relational and context-dependent but diverges from RQM in three important
respects. First, where RQM relativizes facts to observers, Relational Actualism
grounds actualization specifically in causal interaction: a property becomes actual
when an interaction permanently entangles the system with its asymptotic
environment and generates an irreversible causal record. The notion of `observer'
is not doing fundamental work here---every physical interaction that meets this
threshold constitutes an actualization event.

Second, RQM is comfortable with different observers assigning incompatible wave
functions to the same system. The present view holds, by contrast, that there is a
single globally consistent wave function for the universe---following Hartle and
Hawking \citep{HartleHawking1983}---and that apparent perspectival differences
reflect the local availability of information rather than genuine ontological
relativity.

Third, Relational Actualism makes a stronger ontological claim about Hilbert
space: the amplitudes and their interference structure are physically real, not
merely predictive bookkeeping. This claim is supported by the Berry phase,
interference washout experiments, and the results of the quantum eraser discussed
in Section~\ref{sec:experiment}.

% ── Table 1 ──────────────────────────────────────────────────────────────────
\begin{table}[ht]
\centering
\caption{Interpretive landscape: Comparing RAQM with major alternatives.}
\label{tab:landscape}
\small
\begin{tabular}{lllll}
\toprule
\textbf{Framework} & \textbf{Global State} & \textbf{Facts} &
\textbf{Actualization} & \textbf{Randomness}\\
\midrule
Everett \citep{Everett1957}
  & Yes & Relative to branch & None (unitary branching) & Epistemic\\[2pt]
Zurek \citep{Zurek2003}
  & Yes & Left ambiguous & Apparent (einselection) & Epistemic\\[2pt]
GRW/CSL \citep{Bassi2013}
  & Yes & Objective
  & \parbox[t]{3cm}{Stochastic gravitational/\\noise threshold} & Ontic\\[6pt]
Rovelli \citep{Rovelli1996}
  & No & Relative to observer & Perspectival interaction & Ontic\\[2pt]
\textbf{RAQM (This)}
  & Yes & Objective (post-event)
  & \parbox[t]{3cm}{Causal pruning\\(entropy criterion)} & Ontic\\
\bottomrule
\end{tabular}
\end{table}

\subsection{Decoherence and Zurek's Einselection: Engagement and Distinction}

Decoherence theory, developed most systematically by Zurek et al., provides the
dominant technical framework for explaining the quantum-to-classical
transition~\citep{Zurek1981,Zurek2003,JoosZeh1985}. The central idea is that
quantum systems interact with environmental degrees of freedom, causing the reduced
density matrix to rapidly lose its off-diagonal coherences in a preferred
basis---the pointer basis. This process---einselection---explains why macroscopic
objects appear classical.

Decoherence theory is mathematically rigorous and empirically well-confirmed.
However, it faces a residual interpretive problem: decoherence explains why
interference terms become practically unobservable, but it does not explain why a
single outcome occurs. The full density matrix of system-plus-environment remains
in a superposition; only the reduced density matrix loses coherence. Relational
Actualism offers a resolution: the loss of coherence tracked by decoherence
reflects a genuine causal process in which environmental interactions eliminate
possibilities from Hilbert space. Decoherence theory correctly describes the
mathematical structure of this process; the present view adds the ontological
interpretation that possibility-elimination is real and that the resulting
actualized state is genuinely definite, not merely effectively so. The wave
function of the universe is not branching, but pruning.

\subsection{Positioning Against Dynamical Collapse Models}
\label{sec:collapse}

Dynamical collapse models---most prominently the Ghirardi-Rimini-Weber (GRW)
theory, Continuous Spontaneous Localization (CSL), and the gravity-induced
collapse models of Di\'{o}si and Penrose---share RAQM's commitment to objective,
physical collapse~\citep{Bassi2013}. In these frameworks, collapse is a real
physical event, not merely an epistemic update. RAQM is in broad agreement with
this ontological commitment and regards these models as pointing in the right
direction.

However, RAQM diverges from dynamical collapse models in two important respects.
First, the collapse trigger: GRW and CSL invoke a stochastic noise field with a
new fundamental constant (the collapse rate), while Di\'{o}si-Penrose invoke a
gravitational threshold. RAQM instead grounds actualization in the causal
structure of physical interaction itself---specifically, in the irreversible
increase of entanglement entropy with the asymptotic environment produced by an
on-shell interaction. This criterion requires no new fundamental constant and is
basis-independent by construction (defined at the level of the density matrix, not
a perturbative expansion). A full claim of gauge-invariance requires additional
care in theories with constrained Hilbert spaces and edge modes, as discussed in
Section~\ref{sec:criterion}; this is an open refinement.

Second, the domain of application: dynamical collapse models are designed to solve
the measurement problem by modifying the Schr\"{o}dinger equation at a specific
scale. RAQM does not modify the dynamics; it provides an ontological
interpretation of what the existing formalism tracks. The entropy criterion is not
an additional term in the Hamiltonian but a physical condition distinguishing
actualizing from non-actualizing interactions within the existing QFT framework.
For a comprehensive review of dynamical collapse models and their experimental
status, see Bassi et al.\ \citep{Bassi2013}.

% ─────────────────────────────────────────────────────────────────────────────
\section{Mathematical Formalism}
\label{sec:formalism}

\subsection{Pure States, Superposition, and the Born Rule}

A quantum system is described by a state vector $\ket{\psi}$ in a Hilbert space
$\mathcal{H}$. For a two-state system:
\begin{equation}
  \ket{\psi} = \alpha\ket{x_1} + \beta\ket{x_2},\qquad
  \alpha,\beta \in \mathbb{C},\quad |\alpha|^2 + |\beta|^2 = 1.
\end{equation}
The coefficients $\alpha$ and $\beta$ are probability amplitudes. By the Born
Rule, the probability of finding the system in state $\ket{x_i}$ upon measurement
is $|\braket{x_i}{\psi}|^2$. Prior to a context-defining interaction, both terms
in the superposition are physically real in the sense that they contribute to
interference phenomena. The double-slit experiment exemplifies this directly: the
detection probability $|\psi(y)|^2$ contains an interference term
$2\operatorname{Re}(\braket{y}{A}\braket{B}{y})$ responsible for the fringe
pattern, which vanishes exactly when which-path information is available in the
environment.

\subsection{Density Matrices and the Reduced Density Matrix}

The density matrix formalism is essential for describing open quantum systems. For
a pure state $\ket{\psi}$, the density matrix is $\rho = \ket{\psi}\!\bra{\psi}$.
The off-diagonal terms $\alpha\beta^{*}$, $\alpha^{*}\beta$ are the
coherences---the mathematical signature of genuine quantum superposition. Consider
the total system: particle $S$ plus environment $E$. If the environment records
which-path information, the composite state evolves as:
\begin{equation}
  \ket{\Psi_{\mathrm{total}}} = \alpha\ket{x_1}\ket{\varepsilon_1}
  + \beta\ket{x_2}\ket{\varepsilon_2},
\end{equation}
where $\ket{\varepsilon_1}$, $\ket{\varepsilon_2}$ are environmental states
correlated with paths $x_1$, $x_2$. The reduced density matrix of $S$ is obtained
by tracing over $E$:
\begin{equation}
  \rhoS = \operatorname{Tr}_{E}\!\bigl(\ket{\Psi_{\mathrm{total}}}\!
           \bra{\Psi_{\mathrm{total}}}\bigr)
  = |\alpha|^2\ket{x_1}\!\bra{x_1}
  + |\beta|^2\ket{x_2}\!\bra{x_2}
  + \alpha\beta^{*}\braket{\varepsilon_2}{\varepsilon_1}\ket{x_1}\!\bra{x_2}
  + \mathrm{c.c.}
\end{equation}
The off-diagonal terms are suppressed by $\braket{\varepsilon_2}{\varepsilon_1}$---the
overlap between environmental states. When $\braket{\varepsilon_2}{\varepsilon_1}
= 0$, the reduced density matrix becomes diagonal---indistinguishable from a
classical probability distribution. Relational Actualism interprets this as
follows: the environmental interaction has eliminated the non-local path
possibilities from the particle's portion of Hilbert space. The coherences
$\braket{\varepsilon_2}{\varepsilon_1} \to 0$ is the mathematical signature of a
genuine actualization event.

\subsubsection{The Actualization Criterion}
\label{sec:criterion}

For Relational Actualism, the actualization criterion is: \emph{an interaction
constitutes an actualization event when it produces an irreversible increase in
the entanglement entropy of the system with its asymptotic environment}---generating
a robust, redundant causal record recoverable in principle by any future observer
in the causal future of the interaction. This criterion is \emph{sufficient} for
actualization across all regimes treated here. It is also \emph{necessary} in
specific well-controlled models---Markovian open systems and perturbative QFT with
on-shell emission into the far field---where the connection between environmental
orthogonality and irreversible record formation can be made explicit. Whether
necessity holds in full generality is an open problem; a complete proof would
require a factorization-independent, algebraic-QFT formulation of
irreversibility---for instance, via relative entropy monotones on local von Neumann
algebras.

This criterion is basis-independent: entanglement entropy $S(A) = -\operatorname{Tr}(\rho_A
\log \rho_A)$ does not depend on the choice of computational basis for the
system's Hilbert space, and it is defined at the level of the density matrix
rather than the Feynman diagram expansion, making it non-perturbative. A stronger
claim of full gauge-invariance requires additional care. In gauge theories,
Hilbert-space factorization and entanglement depend on the choice of operator
algebra and the treatment of edge modes at the boundary between system and
environment---so `basis-independent' does not automatically entail `gauge-invariant'
in constrained systems. The criterion is gauge-invariant in the sense that it does
not depend on a choice of perturbative gauge fixing or bookkeeping; extending it
rigorously to theories with non-trivial edge mode structure (electrodynamics,
non-Abelian gauge theories) is an open refinement, likely requiring the
algebraic-QFT formulation noted above.

In weakly coupled QFT, on-shell emission into the asymptotic environment is a
\emph{sufficient} condition for actualization: a real, on-shell excitation
satisfying $p^{\mu}p_{\mu} = m^{2}c^{2}$ propagates into the far field, producing
environmental states that are orthogonal (or near-orthogonal) and thereby
generating a lasting, redundant entanglement record. Purely virtual (off-shell)
exchanges---which mediate van der Waals forces, the Casimir effect, and similar
phenomena---typically leave the environment in its ground state and produce only
reversible dynamic shifts (such as the Lamb shift) without generating a persistent
causal record. However, the on-shell condition is \emph{not} asserted to be
necessary: near-field couplings and many-body scattering processes can in
principle produce robust environmental distinguishability without a single
propagating on-shell quantum, and the physically relevant question in all regimes
is whether environmental degrees of freedom end in near-orthogonal states that
persist. The on-shell condition is the clean limiting expression of the
actualization criterion in radiative perturbative regimes, not its general
definition.

This reformulation also handles cases where the on-shell language becomes
ambiguous. In strongly correlated condensed matter systems, a quasiparticle
excitation---a phonon, magnon, or polariton---constitutes an actualization event
when it produces an irreversible entanglement increase in the condensed matter
environment, regardless of whether it maps cleanly onto a fundamental gauge boson.
Similarly, contact interactions and near-field couplings actualize when they
satisfy the entropy criterion, even in the absence of propagating radiation. The
deeper requirement in all regimes is irreversible environmental entanglement.

\paragraph{The Kinematic Snap: the microscopic mechanism of irreversibility.}
The precise physical mechanism that converts reversible potentia into irreversible
actuality can be stated in standard QFT terms without appeal to thermodynamic
avalanches or statistical thresholds. We call this the \emph{Kinematic Snap}.

When two fundamental particles interact, their wavefunctions (potentia) overlap.
At short range, this overlap is mediated by virtual, off-shell bosons satisfying
$p^{\mu}p_{\mu} \neq m^{2}c^{2}$. Virtual bosons are off-shell precisely because
they represent the universe's provisional, reversible exploration of possible
states: they cannot propagate to the asymptotic environment, so they cannot
balance the Local Ledger, and they write no vertex to the DAG. The interaction
remains reversible for as long as only virtual exchange occurs.

Actualization occurs at the kinematic threshold: the precise moment at which the
interaction possesses sufficient energy-momentum to elevate a mediating boson to
the mass shell. At this threshold, the boson transitions to an on-shell real
excitation satisfying $p^{\mu}p_{\mu} = m^{2}c^{2}$. An on-shell boson can
propagate to infinity. The Local Ledger Condition then demands permanent recording:
the energy-momentum transfer is inscribed as an irreversible vertex on the causal
DAG. The continuous, reversible hum of virtual off-shell potentia snaps into a
discrete, permanent causal fact.

\paragraph{Relationship to QFT localization theorems.}
\label{sec:localization}
A natural technical concern is whether identifying actualization with the on-shell
threshold contradicts the Reeh-Schlieder theorem or the failure of Newton-Wigner
localization at sub-Compton scales. It does not, for two reasons.

First, the RA actualization criterion is \emph{not} a sharp spacetime localization
claim. The primary criterion is an irreversible increase in quantum relative
entropy $\Delta S(\rho \| \sigma_0) > 0$ with respect to the vacuum reference
state (Section~\ref{sec:criterion}). The on-shell condition $p^{\mu}p_{\mu} =
m^2 c^2$ is the implementation of this entropic criterion in the perturbative
regime --- it is the signature of an excitation that can propagate to the
asymptotic environment and generate a redundant causal record. It is explicitly
\emph{not} a claim that the interaction is localized at a mathematical point
$x_0^\mu$.

Second, the LSZ reduction formula places asymptotic external states on-shell
at $t \to \pm\infty$. RA uses the on-shell condition as an effective diagnostic
of whether a given exchange has crossed the threshold for environmental
entanglement --- not as a claim that LSZ applies locally at each vertex in a
finite interaction region. The physically relevant object is the $S$-matrix
pole, which is frame-independent and well-defined. What RA adds is the
identification of pole-crossing with the writing of a causal DAG vertex: a
metaphysical claim about what the $S$-matrix pole \emph{means}, not a departure
from the mathematics of interacting QFT.

In strongly coupled regimes where the quasi-particle picture breaks down, the
entropic criterion takes over directly: actualization is the irreversible
increase in $\Delta S(\rho \| \sigma_0)$, and the on-shell language becomes an
approximation. This is an open refinement addressed in RAGC
Derivation~2 (\citep{Sandeman2026RAGC}, Derivation~2),
not a contradiction: the entropic criterion is the fundamental definition,
and the on-shell condition is its exact perturbative implementation.

This mechanism has four immediate consequences:

\begin{enumerate}[leftmargin=*]
  \item \textbf{Irreversibility is kinematic, not statistical.} The threshold is
  the dispersion relation $p^{\mu}p_{\mu} = m^{2}c^{2}$---a binary condition
  grounded in standard QFT, not a probabilistic or thermodynamic criterion. The
  snap is sharp and objective.

  \item \textbf{No observer is required.} An electron and a proton do not require
  a scientist to actualize their binding. Once their interaction produces a real
  on-shell photon carrying away the binding energy ($\sim 13.6$~eV for hydrogen),
  the vertex is written regardless of whether anyone is watching. The mechanism
  is entirely relational and causal.

  \item \textbf{Finite decoherence time is explained.} Decoherence takes a finite
  duration not because of a fuzzy statistical process, but because it takes
  physical time for emitted on-shell bosons to propagate outward and entangle with
  the surrounding environmental degrees of freedom. The decoherence timescale
  $\tau_D \sim \hbar / \Delta E_{\mathrm{env}}$ (Prediction~P2) is precisely the
  light-travel time for on-shell radiation to reach the environmental boundary,
  scaled by the interaction strength.

  \item \textbf{Virtual exchanges are not actualized.} Van der Waals forces, the
  Casimir effect, the Lamb shift---all mediated by off-shell virtual
  exchange---produce no DAG vertices. They are genuine physical effects, but they
  are effects of the potentia, not actualities. This is consistent with the RA
  vacuum prescription: the Casimir energy difference (arising from the actualized
  boundary conditions imposed by the plates) gravitates; the absolute virtual
  zero-point energy does not.

  \item \textbf{Consistency with loop corrections to physical particle masses.}
  A question naturally arises: if virtual exchanges do not actualize, why do
  loop corrections (e.g., pion cloud contributions to the proton mass) shift the
  physical mass of an observable particle? The answer lies in the distinction between
  two types of virtual process.

  \emph{Kind~1 (vacuum sector):} Spontaneous virtual pair creation and annihilation
  with no external particles. These occur in the zero-particle sector, where
  $\mathcal{P}_{\mathrm{act}}|0\rangle = 0$. No actualization ever selects the
  vacuum as its outcome; its energy contributes nothing to the metric source. This
  is the origin of RA's vacuum catastrophe resolution.

  \emph{Kind~2 (particle sector):} Off-shell intermediate states in the propagation
  of a physical external particle. The external particle \emph{does} actualize.
  When it does, it carries its full quantum eigenvalue---including all loop
  contributions from Kind~2 virtuals---into the DAG. These corrections modify
  the energy eigenvalue of the state that actualizes; they are not sources of
  independent DAG vertices.

  The two kinds are separated cleanly by external particle number. Quantum field
  theory conflates them in dimensional regularization, where both appear as formally
  identical loop integrals. RA's physical regulator (the Planck-scale causal
  diamond) makes the separation manifest: a closed vacuum loop with no external
  legs has no actualization event to attach to. A self-energy loop attached to a
  propagating particle's worldline is part of that particle's actualized eigenvalue.
  There is no structural inconsistency between the vacuum catastrophe resolution
  and the physical mass spectrum.
\end{enumerate}

\paragraph{The Unruh vacuum and the stationarity criterion.}
The Kinematic Snap also resolves the Unruh frame-dependence question precisely.
A uniformly accelerated (Rindler) observer detects a stationary thermal bath at
temperature $T_U = \hbar a / (2\pi c k_B)$. This bath is in thermal equilibrium:
its relative entropy with respect to the Rindler vacuum is constant in time
($\Delta S_{\mathrm{rel}} = 0$). No irreversible increase in relative entropy
occurs; no on-shell bosons are being emitted into the asymptotic environment at
a net rate that exceeds absorption. The Rindler thermal bath represents the
universe's stationary potentia as seen by the accelerated observer---it is not
generating new actualization events. Both the inertial and the accelerated
observer therefore agree: in the empty Minkowski vacuum, no actualization is
occurring. The apparent disagreement on particle content is a modal decomposition
artefact, not a disagreement about physical causal facts.

This becomes dramatic at the extreme limit. As acceleration---or equivalently,
by the Einstein Equivalence Principle, gravitational mass-density---approaches
the bandwidth limit $R_{\mathrm{max}}$, the local actualization density of the
DAG saturates. The graph cannot accept further causal links. At this limit, the
kinematic threshold is reached at every interaction simultaneously: the graph
geometrically shears, producing a causal severance. This is the RA account of
what standard GR calls a singularity. There is no infinity; there is a graph cut,
governed by the same Ledger Partition Theorem verified in Lean~4/Mathlib, forcing
the nucleation of a new, causally isolated causal ledger. Recent experiments
measuring the finite duration of quantum transitions~\citep{Minev2019} are
uniquely positioned to probe the kinematic snap mechanism at the laboratory scale,
as discussed in Section~\ref{sec:falsify}.

\paragraph{Frame-independence of the actualization criterion.}
A subtle but important point concerns the Lorentz covariance of the criterion.
If actualization were defined in terms of particle number---``an on-shell particle
is emitted''---it would be frame-dependent: an accelerated (Rindler) observer
detects a thermal bath of particles in the Minkowski vacuum, while an inertial
observer detects none. This would make actualization observer-dependent, which
is precisely what RA rejects.

The entropy criterion resolves this. Relative entropy $S(\rho \| \sigma)$
between states is a diffeomorphism-invariant scalar. A strictly irreversible
increase in relative entropy with respect to the vacuum state is an objective,
coordinate-independent fact. An accelerated and an inertial observer may disagree
on the particle content of an interaction, but they must agree on whether an
irreversible thermodynamic event occurred, because irreversibility is encoded in
the monotone properties of relative entropy---a frame-independent quantity.
The on-shell particle language is the clean limiting expression of this criterion
in weakly coupled, perturbative regimes; the fundamental criterion is the relative
entropy increase, which is covariant in all regimes. The formal proof that
$\mathcal{P}_{\mathrm{act}}$, reformulated in terms of local algebras and modular objects rather than
Fock-space particles, produces a frame-independent set of actualization events is
pursued in the companion paper RAGC~\citep{Sandeman2026RAGC}.

\subsubsection{Lemma: Conservation of Probability under Actualization}

When an actualizing interaction occurs, the continuous pruning of the state space
must strictly conserve total probability. We establish this in two steps,
distinguishing the unconditional update from the post-selection formula.

\textbf{Step 1: The unconditional update.} Let $\{M_k\}$ be the set of Kraus
operators representing the environmental scattering channels, satisfying the
completeness relation $\sum_k M_k^{\dagger}M_k = I$. The unconditional
(pre-selection) state after the interaction is:
\begin{equation}
  \rho_{\mathrm{after}} = \sum_k M_k\,\rho_{\mathrm{before}}\,M_k^{\dagger}.
\end{equation}
Probability conservation follows directly from the completeness relation:
\begin{equation}
  \operatorname{Tr}(\rho_{\mathrm{after}})
  = \operatorname{Tr}\!\Bigl(\sum_k M_k\,\rho_{\mathrm{before}}\,M_k^{\dagger}\Bigr)
  = \operatorname{Tr}\!\Bigl(\rho_{\mathrm{before}}\sum_k M_k^{\dagger}M_k\Bigr)
  = \operatorname{Tr}(\rho_{\mathrm{before}})
  = 1.
\end{equation}
This establishes that the total probability is preserved \emph{across all possible
outcomes}. No probability leaks from the system during the transition from the
possible to the actual.

\textbf{Step 2: The post-selection formula.} Once a specific outcome $k$ is
realized---i.e., once actualization selects a particular environmental
channel---the updated state is renormalized to reflect the realized branch:
\begin{equation}
  \rho_{\mathrm{after}|k}
  = \frac{M_k\,\rho_{\mathrm{before}}\,M_k^{\dagger}}
         {\operatorname{Tr}(M_k\,\rho_{\mathrm{before}}\,M_k^{\dagger})}.
\end{equation}
The denominator $\operatorname{Tr}(M_k\,\rho_{\mathrm{before}}\,M_k^{\dagger}) =
p_k$ is exactly the probability of outcome $k$, as given by the Born Rule. The
post-selection formula geometrically rescales the remaining phase space to the
realized branch, guaranteeing $\operatorname{Tr}(\rho_{\mathrm{after}|k}) = 1$.
The laws of probability remain strictly intact. The unconditional update (Step~1)
establishes conservation across all outcomes; the post-selection formula (Step~2)
gives the state conditional on a specific actualization having occurred.

\subsection{Decoherence Timescales}

The decoherence timescale $\tau_D$ quantifies how rapidly a superposition loses
coherence through environmental interaction. For a particle of mass $m$ in thermal
contact with an environment at temperature $T$:
\begin{equation}
  \tau_D \approx \tau_R \left(\frac{\lambda_{dB}}{\Delta x}\right)^{\!2},
\end{equation}
where $\tau_R$ is the relaxation time, $\lambda_{dB} = \hbar/\sqrt{2mkT}$ is the
thermal de Broglie wavelength, and $\Delta x$ is the spatial separation of the
superposed states. For macroscopic systems, $\Delta x \gg \lambda_{dB}$, and
$\tau_D \ll \tau_R$---decoherence occurs many orders of magnitude faster than
dissipation. For a buckyball (C$_{60}$, mass $\approx 1.2 \times 10^{-24}$\,kg)
at room temperature with path separation $\Delta x \approx 1\,\mu$m, $\tau_D$ is
of order $10^{-19}$\,s.

\subsection{Entanglement Entropy as a Measure of Relationality}

The entanglement entropy provides a quantitative measure of how `relational' a
bipartite system is. For a bipartite state $\rho_{AB}$, the entanglement entropy
of subsystem $A$ is:
\begin{equation}
  S(A) = -\operatorname{Tr}(\rho_A \log \rho_A),
\end{equation}
where $\rho_A = \operatorname{Tr}_B(\rho_{AB})$. $S(A) = 0$ if and only if the
global state is a product state ($A$'s properties are intrinsic). $S(A)$ is
maximized for a maximally entangled state, in which $A$'s properties are entirely
relational---undefined without reference to $B$. This gives precise mathematical
content to the taxonomy of Section~\ref{sec:landscape}: invariant properties
correspond to $S = 0$ for the relevant degree of freedom; intrinsically relational
properties correspond to high-$S(A)$ states. The decoherence process continuously
drives $S(A)$ upward through entanglement with the environment while reducing
local superposition---this is the Hilbert space picture of classicalization.
The formal account of how recursive coarse-graining over the microscopic DAG
generates successive levels of classical description — from chemistry to biology
to agentic intelligence — is developed in the companion paper
RAHC~\citep{Sandeman2026RAHC}.

We can map this process directly onto the geometry of a dynamical phase space. In
Relational Actualism, entanglement entropy functions as the precise quantitative
measure of accumulated causal constraints. Each actualization event physically
constrains the system, binding its future evolution to the state of the
environment. As $S(A)$ increases, the unconstrained volume of the local quantum
state shrinks, funneling the system's trajectory until it is strictly bounded by
its relational network.

\subsection{The Berry Phase as Evidence for the Reality of Hilbert Space}

The Berry phase provides particularly compelling evidence that the geometry of
Hilbert space has physical consequences. When a quantum system's Hamiltonian is
slowly varied along a closed path $C$ in parameter space, the state acquires a
geometric phase:
\begin{equation}
  \gamma_n(C) = i \oint_C \bra{n(\mathbf{R})}\nabla_{\mathbf{R}}\ket{n(\mathbf{R})}
  \cdot d\mathbf{R}.
\end{equation}
This geometric phase depends only on the geometry of the path in parameter
space---specifically, on the solid angle subtended by $C$ in the projective
Hilbert space---not on the rate of traversal or the energy of the state. It has
been experimentally confirmed in systems ranging from photon polarization to
neutron interferometry to topological insulators~\citep{Cohen2019}. The Berry
curvature $\Omega_n = \nabla_{\mathbf{R}} \times \mathbf{A}_n$ (where
$\mathbf{A}_n = i\bra{n}\nabla_{\mathbf{R}}\ket{n}$ is the Berry connection) is
physically real and measurable, strongly supporting the central ontological claim
of Relational Actualism: Hilbert space amplitudes represent real structures whose
geometry constrains and shapes physical processes.

% ─────────────────────────────────────────────────────────────────────────────
\section{Experimental Analysis}
\label{sec:experiment}

\subsection{The Double-Slit Experiment}

The double-slit experiment and its variants provide the central experimental
scaffold for the present argument. In the scenario with one open slit, the pattern
on the detection screen is consistent with point-like entities. With two slits, a
wave interference pattern appears. In the density matrix framework, this is
precisely the situation of non-zero off-diagonal coherences:
$\braket{\varepsilon_2}{\varepsilon_1} \approx 1$, because no which-path
information has been recorded in the environment. When a detector is placed at one
slit, the environment acquires which-path information,
$\braket{\varepsilon_2}{\varepsilon_1} \to 0$, and the interference pattern
vanishes. Whether the detector `fires' or not is irrelevant---what matters is
whether the environmental entanglement entropy has increased irreversibly.

\subsection{The Buckyball Experiment and the Washout Effect}

The buckyball (C$_{60}$) interference experiment of Hackerm\"{u}ller et
al.~\citep{Hackermuller2004} is particularly illuminating because it demonstrates
that the transition from quantum to classical behavior is continuous, not
discontinuous, and is governed by information content. At low temperatures,
buckyballs exhibit clear wave interference through a nanograting. As temperature
increases, the buckyballs emit thermal radiation---photons whose wavelength encodes
which-path information. At a critical temperature the interference pattern washes
out entirely. A variation by Hornberger et al.~\citep{Hornberger2003} achieved the
same washout by degrading the vacuum.

The key physical quantity is the overlap $\braket{\varepsilon_2}{\varepsilon_1}$
between environmental states correlated with the possible paths. The washout is
continuous and quantitative, matching the gradual reduction of off-diagonal
elements in the reduced density matrix. The energy threshold $E_{\mathrm{wash}}
\approx hc/d$ (for slit separation $d \approx 100$\,nm gives $E_{\mathrm{wash}}
\approx 12.4$\,eV) is consistent with experimental results.

\subsection{Delayed-Choice and Quantum Eraser Experiments}

The delayed-choice experiment (Wheeler) demonstrates that the presence or absence
of an interference pattern is determined by the total informational state of the
system at the moment of detection---the relevant quantity is
$\braket{\varepsilon_2}{\varepsilon_1}$ evaluated over the full system history.
This explains why there is no retrocausal signaling.

The quantum eraser experiment~\citep{ScullyDruhl1982} completes the picture: by
destroying which-path information (via a second polarizer), coherence is
restored---the off-diagonal terms re-emerge in the reduced density matrix of the
signal photons. This demonstrates bidirectionality: actualization can be `undone'
if the environmental record is erased \emph{before} a macroscopic thermodynamic
record is irreversibly established. This is critical: the quantum eraser does not
reverse an actualization; it reveals that no actualization had yet occurred,
because the thermodynamic horizon had not been crossed. This is strong evidence
that irreversible environmental entanglement---not the physical presence of a
detector---is the operative quantity.

\subsection{Quantum Maze: Entanglement and Nonlocal Constraints}

The quantum maze experiment uses parametric down-conversion to produce entangled
photon pairs. Because the pairs are entangled, they constitute a single quantum
system; a measurement on the idler photon instantaneously constrains the state of
the signal photon regardless of spatial separation. The present framework
interprets this as follows: the entangled pair has a joint wave function that is
non-separable. When the idler is detected, the actualization criterion is satisfied
and the signal photon's reduced density matrix loses its off-diagonal terms. This
nonlocality is not a violation of relativity because no information is transmitted
faster than light; both photons belong to a single nonlocal quantum system whose
wave function is defined in the configuration space of both particles jointly.

% ─────────────────────────────────────────────────────────────────────────────
\section{Discussion}
\label{sec:discussion}

\subsection{The Many-Worlds Rebuttal}

The Everettian Many-Worlds Interpretation
(MWI)~\citep{Everett1957,SaundersEtAl2010} holds that the universal wave function
never collapses and evolves unitarily always. MWI has notable virtues: it takes
the formalism of quantum mechanics at face value and requires no additional
collapse postulate. However, it faces several serious difficulties.

First, the preferred basis problem. While einselection provides a basis-selection
mechanism via decoherence, branching must occur in some basis. The pointer basis
selected by decoherence depends on the structure of the environment, which must be
specified. The branches of MWI are not fundamental features of the wave function
but artifacts of a particular factorization of Hilbert space.

Second, the probability problem. Deriving the Born rule probabilities within MWI
is notoriously difficult. If all branches occur, in what sense are some outcomes
more probable than others? Deutsch-Wallace decision-theoretic arguments and
branch-counting approaches both face objections~\citep{SaundersEtAl2010}.

Third, and most fundamentally: MWI denies that anything actually happens. No
branch is more real than any other; the appearance of a definite outcome is an
illusion relative to our branch. This conflicts with the most basic datum of
experience---that something, specifically, occurs---and with the scientific
methodology that requires comparing theory with a definite observational record.
Relational Actualism preserves the ontological primacy of actual events: something
happens, with probability given by the Born rule, and this happening dynamically
eliminates specific possibilities. The wave function prunes rather than branches.

Fourth, the ontological parsimony argument cuts both ways. MWI is often praised
for positing only one entity (the wave function) with no collapse. But it
proliferates universes---an ontological cost not obviously lower than positing a
single universe with genuine randomness.

\subsection{Ontic Randomness, the Equal-Rate Limit, and the Kinematic Snap}
\label{sec:ontic}

The framework developed here implies that the determinism/indeterminism distinction
is not fundamental but contextual. Nature is deterministic when the causal and
logical context has reduced the possibilities for a variable to exactly one; it is
irreducibly random when multiple possibilities remain viable.

It is crucial to note that while classical chaos in phase space represents purely
epistemic uncertainty about a fully determined trajectory, the randomness in RAQM
is \emph{ontic}: the trajectory itself is physically undetermined until
actualization dynamically prunes the space. This distinction between epistemic and
ontic randomness is load-bearing for the framework and must not be elided.

\paragraph{The Kinematic Snap grounds ontic randomness precisely.}
The competing Poisson process derivation of the Born rule (Section~\ref{sec:criterion})
makes the locus of randomness explicit. Each volume element $d^3x$ in the wavefunction
support independently races to cross the kinematic threshold $p^\mu p_\mu = m^2c^2$.
The time at which any given volume element crosses the threshold is governed by the
quantum amplitude for the mediating boson to go on-shell --- itself an underdetermined
quantity until it occurs. This is not ignorance of a hidden variable. There is no
fact of the matter, in the prior causal history, about when the threshold will be
crossed at any specific location. The snap time at each location is ontic, not
epistemic.

\paragraph{The equal-rate limiting case.}
A natural question arises: if two locations have exactly equal virtual exchange
rates --- as in a perfectly symmetric double-slit experiment or an equal-amplitude
two-state superposition --- is there a fact of the matter about which location fires
first? And if not, has determinism been smuggled back in through the back door?

The answer is no, and the argument is precise. In the equal-rate case, the
competing Poisson process assigns $P(x_1) = P(x_2) = 1/2$. Neither location is
favoured by the prior causal history. The snap at $x_1$ versus $x_2$ is not
determined by any prior fact --- it is a genuinely new, uncaused addition to the
causal DAG. This is not a failure of the framework; it is the framework's central
claim. The actualization event at $x_0$ is ontologically primitive: the transition
from possible to actual is the irreducible locus of contingency in the universe.

To ask ``but why $x_0$ rather than $x_1$?'' is to ask for a prior cause of an
event that, by definition, is the first link in a new causal chain. The question
presupposes that every event must have a prior determining cause --- which is
precisely the deterministic assumption that ontic randomness rejects. In RA, the
actualization event is not caused by prior events in the sense of being logically
entailed by them. Prior events determine only the probability distribution over
possible actualization locations. Which possibility becomes actual is the
universe's irreducible freedom --- the only place where the possible genuinely
outstrips the determined.

\paragraph{Is perfect symmetry physically realizable?}
In practice, no experimental setup achieves perfect equality of virtual exchange
rates at two locations. Environmental asymmetries, vacuum fluctuations, and the
dense causal history of the laboratory break the symmetry at some level. This is
not a loophole that restores determinism. Even with arbitrarily small asymmetry
$\epsilon$, the outcome at the less-favoured location occurs with probability
$1/2 - \epsilon/2$ --- which approaches $1/2$ but never reaches zero. The
randomness is not resolved by the asymmetry; it is merely shifted. The fundamental
ontic randomness of the snap survives in every physically realizable case.

\paragraph{Comparison: why classical chaos is different.}
Classical chaos is deterministic at the level of individual trajectories. Given
perfect initial conditions, the trajectory is fixed; only our ignorance of those
conditions produces apparent randomness. In the RA Kinematic Snap, there are no
hidden initial conditions that would determine the snap location if we knew them.
The quantum amplitudes governing the threshold crossing time at each location are
not waiting to be discovered --- they are underdetermined until the snap occurs.
The randomness is in the world, not in our description of it.

This is the distinction Bohr was groping toward but could not state precisely
without the ontological framework RA provides. The wave function is not a
description of our ignorance; it is a description of the universe's genuine
openness. Actualization is the moment that openness closes.

Classical behavior is the dense-graph regime of the same discrete dynamics:
as $\tau_D \to 0$ and $\braket{\varepsilon_2}{\varepsilon_1} \to 0$,
the Poisson-CSG parameter $\mu \gg 1$ and the graph is so densely
actualized that quantum indeterminacy has been eliminated at every scale
of interest. Large, warm, complex systems are classical not because they
are fundamentally different from quantum systems but because their causal
graph is far above the Causal Firewall threshold. This resolves the
Heisenberg cut without invoking a cut at all: there is no boundary between
quantum and classical regimes, only a gradient in $\mu(x)$.%
\footnote{Schr\"{o}dinger's cat deserves a brief defense---of the cat. The genuine
locus of quantum indeterminacy in the thought experiment is the isotope: a small,
isolated system whose decay is irreducibly random in exactly the sense RAQM
identifies. A real cat is continuously actualizing---metabolizing, moving, emitting
heat---a dense internal network of on-shell interactions that pins it firmly in the
classical limit regardless of what the isotope is doing. RAQM rescues the world we
actually see: the cat remains alive and indignant in its box until the decay event
propagates through the mechanism and ends the matter. The paradox dissolves not
because observation collapses a superposition, but because the cat was never in
one.}

\subsection{The Global Wave Function, Relativistic Causality, and Microcausality}
\label{sec:relativity}

There is one wave function for the
universe~\citep{HartleHawking1983}. Every interaction updates it. However,
positing a global update to the wave function immediately invites a relativistic
critique: a simultaneous global update seemingly requires a preferred universal
`now,' which would violate Lorentz invariance. This tension is resolved by
formalizing the global wave function using the Tomonaga-Schwinger equation. Rather
than evolving according to a single global time parameter $t$, the universal
quantum state is defined as a functional $\Psi[\sigma]$ on an arbitrary spacelike
hypersurface $\sigma$:
\begin{equation}
  i\hbar\,\frac{\delta\Psi[\sigma]}{\delta\sigma(x)} = H(x)\,\Psi[\sigma].
\end{equation}
An actualizing interaction updates the state functional along any valid spacelike
hypersurface passing through $x$; the sequence of actualization events remains
globally consistent regardless of how an observer slices spacetime.

The prohibition on superluminal signaling is established rigorously through
algebraic microcausality. For any two spacelike-separated regions $A$ and $B$,
where $(x_A - x_B)^2 < 0$:
\begin{equation}
  [\hat{O}_A,\,\hat{O}_B] = 0.
\end{equation}
An actualization event at $A$ instantly restricts the joint probability
distribution---pruning possibilities within the global configuration space---but
microcausality ensures that the \emph{non-selective} (outcome-averaged)
actualization map cannot causally force a measurable change in the expectation
values of observables at spacelike-separated $B$. For a \emph{conditioned}
(post-selected) branch, the remote conditional state can in general differ from the
pre-actualization state; however, because the specific outcome at $A$ is not
controllable from $B$, this conditional difference cannot be exploited to transmit
information superluminally. There are genuine nonlocal constraints that are
enforced whenever results are compared, a process that requires classical signaling
at sub-luminal speeds.

A further complementary bound for non-relativistic systems: the Lieb-Robinson
bound~\citep{LiebRobinson1972} establishes that in quantum lattice systems, the
speed of propagation of correlations is bounded by an effective `speed of light'
determined by local interaction strength. This points toward a natural extension of
the present framework to non-relativistic quantum systems---a direction pursued in
the companion paper (Section~\ref{sec:future}), where the Causal Firewall defines
the non-relativistic limit in which the local actualization history is universally
agreed upon.

\subsubsection{Consistency Condition for Actualization Events}
\label{sec:consistency}

The Tomonaga-Schwinger equation guarantees foliation-independence for unitary
evolution: for any two spacelike-separated points $x$ and $y$,
\begin{equation}
  \left[\frac{\delta}{\delta\sigma(x)},\,
        \frac{\delta}{\delta\sigma(y)}\right]\Psi[\sigma] = 0,
  \qquad (x-y)^2 < 0,
\end{equation}
follows from microcausality. In free field theory this is the statement $[H(x),
H(y)] = 0$ for $(x-y)^2 < 0$; in interacting, renormalized QFT the Hamiltonian
density $H(x)$ must be understood as an operator-valued distribution whose local
commutativity is one of the Wightman/Haag-Kastler axioms of algebraic
QFT~\citep{HaagKastler1964}, not an automatic consequence of naive density
expressions. For the purposes of RAQM, we work within the framework of locally
covariant QFT where local commutativity is assumed as an axiom; a fully
constructive proof in specific interacting models (e.g., an Unruh-DeWitt detector
coupled to a free scalar field) can be made explicit and would constitute a
valuable concrete verification. What requires independent argument beyond this
unitary foundation is whether RAQM's actualization map---the non-unitary operation
that selects an outcome and updates $\Psi[\sigma]$ at an interaction vertex---also
satisfies an appropriate consistency condition.

Let an actualization event occur at spacetime point $x$, and let $\sigma_1$ and
$\sigma_2$ be any two spacelike hypersurfaces both passing through $x$, related by
a local deformation in the causal future of $x$. The RAQM consistency condition
requires that the probability distributions over all observable outcomes at any
point $y$ with $(x-y)^2 < 0$ be identical whether computed from $\Psi[\sigma_1]$
or $\Psi[\sigma_2]$ following actualization at $x$.

The argument proceeds in two steps. First, unitary evolution prior to actualization:
the standard T-S consistency condition guarantees that $\Psi[\sigma_1]$ and
$\Psi[\sigma_2]$ encode identical physical content up to and including the
actualization vertex at $x$. Second, the actualization map: RAQM's actualization
is grounded in a local interaction at $x$ satisfying the irreversible entanglement
criterion, and therefore acts exclusively on the degrees of freedom in a
neighbourhood of $x$. For any observable $\hat{O}_B$ supported in a region $B$
spacelike-separated from $x$, algebraic microcausality gives $[\hat{O}_A,
\hat{O}_B] = 0$, where $\hat{O}_A$ is any operator supported near $x$. It follows
that the \emph{non-selective} (outcome-averaged) actualization map---formally
$\sum_k \inst_k$ where $\{\inst_k\}$ is the localized quantum
instrument---leaves the expectation value of every observable at
spacelike-separated $B$ unchanged:
\begin{equation}
  \operatorname{Tr}\!\left[\hat{O}_B \left(\sum_k \inst_k\right)\!(\rho)\right]
  = \operatorname{Tr}[\hat{O}_B\,\rho].
  \label{eq:nosig}
\end{equation}
This is the rigorous no-signaling statement. It should be noted that for a
\emph{conditioned} (post-selected) outcome $k$, the remote conditional state
$\rho_{B|k}$ can in principle differ from the pre-actualization $\rho_B$;
no-signaling is preserved because the outcome $k$ is not locally controllable from
$B$. The foliation-invariance claim applies strictly to the non-selective map and
to the resulting observable statistics, not to individual conditioned branches.

One further clarification. This argument establishes that observable statistics at
spacelike-separated points are foliation-independent: no measurement at $B$ can
distinguish which hypersurface was used to carry the actualization at $x$. It does
not establish that the global state functional $\Psi[\sigma]$ is literally
identical on $\sigma_1$ and $\sigma_2$ following actualization---since actualization
is a non-unitary operation, the global descriptions will in general differ. What
is invariant is the empirical content: the full set of outcome probabilities
accessible to any observer in the causal future of an actualization event. This is
the physically relevant notion of foliation-independence for RAQM.

\subsubsection{Foliation-Independence of the Kinematic Snap}
\label{sec:foliation-snap}

The Kinematic Snap mechanism raises a sharpened version of the
foliation-independence question. The competing Poisson process runs simultaneously
across the entire wavefunction support: every volume element $d^3x$ where
$|\psi(\mathbf{x})|^2 > 0$ is independently racing to cross the kinematic
threshold. Different spacelike hypersurfaces will slice this distributed,
simultaneous process differently. We must verify that the snap location and the
Born probabilities are hypersurface-independent.

\paragraph{The snap event is a spacetime point, not a hypersurface event.}
The first and most important observation is that the Kinematic Snap, when it
occurs, is a local event at a specific spacetime point $x_0^\mu$. A virtual
exchange at $x_0^\mu$ crosses the threshold $p^\mu p_\mu = m^2c^2$ and an
on-shell boson is emitted from $x_0^\mu$. This is a pointlike event in
spacetime --- it has a definite location in all reference frames, because spacetime
points are covariant. Different observers may assign different spatial coordinates
and different times to $x_0^\mu$, but they agree that the snap occurred there.

The snap is therefore not ``which location fires first on a given hypersurface''
but rather ``which spacetime point is the first to emit an on-shell boson.'' The
ordering is causal, not hypersurface-dependent.

\paragraph{The competing Poisson process in covariant form.}
The snap rate at spacetime point $x^\mu$ is proportional to the local probability
density in the rest frame of the interaction. For a Dirac field with conserved
four-current $j^\mu = \bar{\psi}\gamma^\mu\psi$, the timelike component
$j^0 = \psi^\dagger\psi = |\psi(x)|^2$ gives the probability density in the local
rest frame. The Lorentz-covariant snap rate is
\begin{equation}
  r(x^\mu) \;\propto\; j^\mu(x)\,u_\mu(x),
\end{equation}
where $u^\mu(x)$ is the future-directed unit timelike vector of the local rest
frame. In the rest frame, $u^\mu = (1,\mathbf{0})$ and $r \propto j^0 = |\psi(x)|^2$.
Note that the contracted scalar $j^\mu j_\mu$ is a Lorentz invariant but does not
equal $|\psi|^2$ in general; the physically relevant quantity is the rest-frame
probability density $j^\mu u_\mu$.

The competing process is therefore between spacetime volume elements $d^4x$, not
between spatial volume elements $d^3x$ on a preferred hypersurface. The
probability that the snap occurs in $d^4x$ at $x^\mu$ before it occurs anywhere
else in the support of the wavefunction is:
\begin{equation}
  dP(\text{snap at } x^\mu) = \frac{r(x^\mu)\,d^4x}{\int_{\mathrm{supp}(\psi)} r(x'^\mu)\,d^4x'},
\end{equation}
where the denominator integrates over the spacetime support of $\psi$, ensuring
the normalisation is finite and well-defined. The snap location is determined by
the four-dimensional local probability density, not by a preferred time coordinate.

\paragraph{Born rule is foliation-independent.}
For a non-relativistic wavepacket localized in time (a pulse of finite duration),
the integration over $dt$ in the denominator is the same for all observers up to
time-dilation factors that cancel in the ratio. For each spatial location, the
probability of snap at $\mathbf{x}$ integrates to $|\psi(\mathbf{x})|^2$
regardless of which spacelike hypersurface is used to evaluate it, because the
time-integrated rate at each spatial location is proportional to $|\psi(\mathbf{x})|^2$
in every frame. The Born rule is therefore foliation-independent.

\paragraph{The snap time is frame-dependent; the snap event is not.}
Different observers will assign different times $t_0$ to the snap event at
$x_0^\mu$. An observer moving relative to the apparatus will see the snap as
occurring earlier or later. But they will agree on: (a) that the snap occurred at
$x_0^\mu$; (b) that the Born probability for $x_0^\mu$ is $|\psi(x_0)|^2$; and
(c) that the wavefunction evolves from $x_0^\mu$ onward. The snap time is
frame-dependent in the same sense that any spacetime coordinate is
frame-dependent. The physical content of the snap --- the vertex written to the
DAG, the on-shell emission, the irreversible entanglement record --- is a
covariant fact.

\paragraph{What remains open.}
The argument above establishes foliation-independence for the snap location and
Born probabilities in perturbative QFT where the competing Poisson process
approximation is valid. Two refinements require further work:

\begin{enumerate}[leftmargin=*]
  \item \textbf{Entangled multi-particle snaps.} When two entangled particles
  interact with separate detectors at spacelike separation, the snap at detector 1
  and the snap at detector 2 must be shown to be consistent regardless of which
  snap is described as occurring ``first'' in any given frame. The no-signaling
  result (Section~\ref{sec:consistency}) guarantees that observable statistics are
  frame-independent, but the joint snap probability for entangled pairs requires a
  careful treatment using the covariant four-current density for the joint system
  rather than individual single-particle currents.

  \item \textbf{Non-perturbative regime.} The Poisson process approximation
  assumes the snap events are rare and independent. In strongly coupled regimes,
  the snap events may not be statistically independent, and the covariant rate
  formula requires the full AQFT formulation. The finite-dimensional version of
  the required frame-independence result has now been formally verified in
  Lean~4: the theorem \texttt{frame\_independence} in
  \texttt{RA\_AQFT\_Proofs\_v10.lean} establishes that
  $S(U\rho U^\dagger \| \sigma_0) = S(\rho \| \sigma_0)$ for any unitary $U$,
  using the unitary-invariance of relative entropy together with
  Poincar\'{e}-invariance of the vacuum as a physics axiom. The proof uses the
  \texttt{Matrix.cfc\_conj\_unitary} lemma from the Lean-QuantumInfo
  library~\citep{Meiburg2025LQI}. The remaining target is the
  infinite-dimensional extension: showing that
  $\mathcal{P}_{\mathrm{act}}$, defined via local algebras, commutes with wedge
  algebra modular flow in the Tomita-Takesaki sense.
\end{enumerate}

The second open item (non-perturbative regime / full AQFT extension) has been
partially resolved: the frame-independence claim is now a proved theorem in finite
dimensions. The physical picture is established and the finite-dimensional
Lean verification is complete; the remaining target is the infinite-dimensional
algebraic QFT extension to type~III$_1$ von Neumann algebras. A collaborator
with expertise in Tomita-Takesaki modular theory would be well-positioned to
complete this step.

\subsection{Free Will, Creativity, and the Constrained Observer}

A provocative implication of the framework deserves careful treatment. RAQM
forecloses libertarian free will as causally incoherent. An uncaused cause would
be a causal contradiction in a universe whose consistency is guaranteed by the
relational structure of actualization. Every actualizable property is intrinsically
relational: it has no determinate value except in the context of a specific causal
interaction. A choice that was causally independent of all prior context would be a
property with a determinate value independent of all relational context---which is
exactly what RAQM rules out as ontologically incoherent.

What RAQM does preserve is genuine creativity. The universe described by RAQM is
not a clockwork in which all outcomes are fixed in advance. At the quantum level,
actualization outcomes are irreducibly random when multiple possibilities remain
viable. This randomness is not mere noise---it is the universe's genuine openness
to novelty. Each new actualization event adds to the causal record of the universe
in ways that were not fixed by prior history. Freedom, understood as the genuine
openness of the future and the real causal efficacy of minds in shaping which
possibilities are actualized, is not merely compatible with RAQM---it is demanded
by it.

\subsection{Empirical Support and Falsifiability}
\label{sec:falsify}

A foundational critique often leveled at ontological interpretations of quantum
mechanics is the problem of empirical equivalence. Relational Actualism goes beyond
empirical equivalence by positing that actualization is a continuous, physical
process of possibility-elimination occurring over a finite duration, driven by
causal interactions, rather than an instantaneous mathematical discontinuity.

This prediction is actively supported by recent experimental breakthroughs.
Measurements of artificial atoms in superconducting circuits have demonstrated
that quantum jumps are not instantaneous, discontinuous events, but continuous and
even reversible transitions occurring over finite timescales~\citep{Minev2019}.
Similarly, attosecond chronoscopy of photoemission~\citep{Pazourek2015} has shown
that the duration of an electron's transition is strictly governed by the local
structural and interaction environment. These results align with RAQM's mechanism,
where actualization is governed by the continuous decay of off-diagonal coherences
($\braket{\varepsilon_2}{\varepsilon_1} \to 0$) and bounded by local interaction
strengths.

By anchoring the transition of the possible to the actual in these measurable,
finite dynamics, Relational Actualism distinguishes itself empirically from
instantaneous collapse models. We now state the specific predictions more
precisely.

\paragraph{Prediction 1: Finite actualization propagation velocity.}
RAQM predicts that the actualization front --- the boundary between the superposed
and actualized regions of a many-body system --- propagates at a finite velocity
bounded by the local interaction strength, not instantaneously. In a quantum
lattice system with nearest-neighbor coupling $J$, the Lieb-Robinson
bound~\citep{LiebRobinson1972} establishes that correlations propagate at most at
velocity $v_{\mathrm{LR}} \sim J \cdot \ell / \hbar$, where $\ell$ is the lattice
spacing. RAQM predicts that the actualization front propagates at or below this
bound. The testable deviation from instantaneous collapse models: in a spatially
extended many-body system, the decoherence of a distant site $B$ following an
actualization event at site $A$ should exhibit a time delay $\Delta t \geq |r_A -
r_B| / v_{\mathrm{LR}}$ before the off-diagonal coherences at $B$ begin to
collapse. Instantaneous collapse models predict no such delay.

\paragraph{Prediction 2: Environment-dependence of transition duration.}
RAQM predicts that the duration of a quantum transition is not an intrinsic
property of the system alone, but is governed by the local environmental coupling
strength. Specifically, the actualization timescale $\tau_{\mathrm{act}}$ should
scale as $\tau_{\mathrm{act}} \sim \hbar / \Delta E_{\mathrm{env}}$, where
$\Delta E_{\mathrm{env}}$ is the characteristic energy scale of the environmental
interaction driving the entanglement. This is consistent with the Pazourek et al.\
attosecond chronoscopy results~\citep{Pazourek2015}, where the ionization duration
depends on the local Coulomb environment, and with Minev et al.'s demonstration
that the jump duration depends on the coupling to the measurement apparatus. The
prediction is that systematically varying the environmental coupling in a controlled
superconducting circuit experiment should produce a proportional variation in
transition duration --- a clean, quantitative test.

\paragraph{Prediction 3: Reversibility as a diagnostic of actualization threshold.}
RAQM distinguishes pre-actualization dynamics (reversible, because the
thermodynamic horizon has not been crossed) from post-actualization dynamics
(irreversible, because a permanent causal record has been established). The quantum
eraser experiments~\citep{ScullyDruhl1982} and the Minev et al.\ jump-reversal
results~\citep{Minev2019} both demonstrate that what appears to be a collapse can
be reversed if the environmental record is erased before thermodynamic
irreversibility is established. RAQM predicts a sharp threshold between these two
regimes, determined by the environmental entanglement entropy crossing a critical
value. Below threshold: the transition is reversible in principle. Above threshold:
the causal record is permanent and the transition cannot be reversed regardless of
intervention. Identifying this threshold experimentally --- the precise point at
which reversal becomes impossible --- would constitute a direct measurement of the
actualization criterion.

\paragraph{Distinguishing RAQM from competing interpretations.}
All three predictions above are shared by RAQM and \emph{not} by the following
alternatives: (a) instantaneous collapse models (GRW, CSL) predict no propagation
delay and no environment-dependence of the collapse threshold beyond the stochastic
rate parameter; (b) Everettian Many-Worlds predicts no sharp threshold and no
genuine irreversibility --- only apparent irreversibility from within a branch; (c)
decoherence-only models predict environment-dependence of decoherence timescales
but deny that a sharp actualization threshold exists. RAQM's predictions are
therefore genuinely discriminating, not merely reframings of existing results.

\subsection{Future Directions: Curved Spacetime and the Arrow of Time}
\label{sec:future}

While the present framework resolves the measurement problem within the domain of
flat Minkowski spacetime, grounding actualization in irreversible environmental
entanglement naturally invites extension into non-inertial reference frames and
curved spacetime. Because the definition of an on-shell excitation is
frame-dependent via the Unruh and Hawking effects, RAQM implies that the rate and
physical reality of wave function pruning must be inherently frame-dependent,
deeply tying the structure of actualization to the causal horizons of the observer.

Furthermore, by defining actualization as the irreversible elimination of
possibilities, this framework provides a purely ontological grounding for the arrow
of time. The formal extension of Relational Actualism to causal severance at event
horizons and non-inertial thermodynamics is developed in the companion paper
\emph{Relational Actualism Gravity and Cosmology} (RAGC)~\citep{Sandeman2026RAGC}.
That paper proves two results that flow directly from the actualization criterion
established in Section~\ref{sec:criterion} of the present work: first, that the
RA-modified stress-energy operator has zero vacuum expectation value, structurally
eliminating the vacuum energy contribution to the metric source term; and second,
that for any causal severance of the spacetime DAG, the local conservation
condition holds independently on each disconnected subgraph, formally establishing
the energy bookkeeping dissolution across event horizons. The second result---the
Ledger Partition Theorem---has been \textbf{formally verified in Lean~4/Mathlib},
establishing the logical soundness of its discrete graph-theoretic core.

RAGC further makes three empirical predictions that follow structurally from the
actualization criterion established here: that the inferred value of $H_0$
correlates with the mean actualization density along the line of sight (providing
a geometric mechanism for the Hubble Tension without dark energy); that the vacuum
energy contribution to the metric is structurally zero rather than fine-tuned; and
that daughter spacetimes nucleated at causal severances inherit an effective
cosmological constant determined by the local actualization flux at nucleation.

The companion paper \emph{Relational Actualism Complexity and Intelligence}
(RACI)~\citep{Sandeman2026RACI} extends the framework to biological complexity and
agency, proving that the Markov blanket isomorphism and the thermodynamic steering
cost bound are structurally grounded in the same actualization criterion, and
making four empirical predictions including a physically derived constraint on the
distribution of life in the universe. Core theorems in RACI have been formally
verified in Lean~4/Mathlib. The load-bearing Condition C1$'$ of the Assembly Index
Correspondence theorem---that each JOIN operation in molecular assembly corresponds
to at least one irreversible actualization event (a coarse-grained lower bound
robust to condensed-phase polymer chemistry)---has been established by a variational
argument (the variational principle guarantees all stable covalent bond formation is
exothermic, mandating irreversible energy dissipation to satisfy the Local Ledger
Condition) and confirmed by B3LYP/6-311+G* density functional calculations across
the five primary bond types of biological assembly (C-C, C-N, C-O, S-S, P-O), all
verified exothermic. The Assembly Index Correspondence theorem therefore holds
across the full chemical alphabet of biological complexity.

% ─────────────────────────────────────────────────────────────────────────────
\section{Conclusion}
\label{sec:conclusion}

Einstein put himself on the path to significant discoveries by accepting that the
speed of light is a constant of Nature and working out the consequences. What he
could not accept was a universe where true randomness is possible. Yet this is what
nature really seems to be like. Relational Actualism accepts indeterminacy and
works out its consequences.

A companion paper, \emph{Relational Actualism: The Topology of Matter}
(RATM)~\citep{Sandeman2026RATM}, provides Lean-verified proofs that the
BDG causal topology types in 4D correspond exactly to the SM particle
spectrum: the sequential/topological causal past partition equals the
lepton/quark distinction, and colour confinement follows as a theorem
with finite confinement lengths $L=3$ (gluon) and $L=4$ (quark).

A further recent result demonstrates the full scope of the programme:
the strong coupling constant is now proved directly from the BDG path weight
at the self-dual point~\citep{Sandeman2026BDG}:
\begin{equation}
  \alpha_s(m_Z) = \frac{1}{\sqrt{c_2 c_4}} = \frac{1}{\sqrt{72}} = 0.11785
  \quad (\text{PDG: } 0.11800,\; 0.13\%),
  \label{eq:alphas_raqm}
\end{equation}
using only the four BDG integers and the same Alexandrov causal-diamond
geometry that underpins the actualization criterion central to this paper.
Together with the Wyler formula for $\alpha_{\mathrm{EM}}$, both Standard
Model coupling constants --- and the dark matter/baryon ratio
$\Omega_{\mathrm{CDM}}/\Omega_b = 5.40$ (Planck~2018: $5.416$) --- follow
from the integers that encode why spacetime is four-dimensional.
The same ontological commitment that dissolves the measurement problem
also constrains the particle physics of the universe.



Properties are of three kinds: invariant (always actualized), classical (actualized
deterministically through interaction), and intrinsically relational (actualized
only in specific contexts of interaction, with genuine ontic randomness when more
than one possibility remains). Crucially, accepting this intrinsic randomness does
not require abandoning relativistic causality.

By formalizing the global wave function using the Tomonaga-Schwinger equation
within the Hartle-Hawking framework, the wave function of the universe is real,
dwelling in a Hilbert space whose geometry (as evidenced by the Berry phase) is
physically meaningful. Possibilities and actualities are both real; manifest
reality is the constant process of the possible becoming actual through
interaction---a dynamic breathing of the local phase space volume, where degrees of
freedom are physically pruned and, prior to thermodynamic irreversibility, can be
restored.

Because this global actualization is foliation-independent---as established by a
derived consistency condition showing that RAQM's actualization map leaves
observable statistics at spacelike-separated points invariant under hypersurface
choice (equation~\ref{eq:nosig}), applying rigorously to the non-selective
map---and rigorously bounded by microcausality, it enforces nonlocal logical
constraints in configuration space without ever permitting superluminal signaling
in Minkowski spacetime.

This framework distinguishes itself from Relational Quantum Mechanics by grounding
actualization in causal interaction rather than perspectival observation. It
distinguishes itself from dynamical collapse models by grounding the collapse
trigger in irreversible environmental entanglement rather than a stochastic
threshold or new fundamental constant. It distinguishes itself from
decoherence-only approaches by adding the ontological claim that
possibility-elimination is genuine, not merely apparent. And it distinguishes
itself from Many-Worlds by preserving the primacy of actual events and the reality
of a single history.

The universe may not always know what it is going to do, but it always knows what
it has done---and this non-locally constrains what is then possible, progressively
funneling its unconstrained degrees of freedom into a definitive, classical record.
There is a simple, if profound, logic to all of it. One would think that Einstein
would have liked that.

% ─────────────────────────────────────────────────────────────────────────────
\section*{Acknowledgements}
The author acknowledges the use of Anthropic's Claude and Google's Gemini during
the preparation of this manuscript. These AI tools were utilized as conversational
assistants to help with structural editing, prose refinement, brainstorming
conceptual analogies, and generating the overall formatting of the manuscript. The
author takes full and sole responsibility for the original theoretical framework,
the accuracy of the physical arguments, and the final intellectual content of the
paper.

\smallskip\noindent
\textbf{Ethics statement.}
This article does not contain any studies with human participants
or animals performed by the author.

\section*{Declarations}
\textbf{Funding:} Not applicable.\quad
\textbf{Competing interests:} The author declares no competing interests.

\textbf{Preprint availability:} This manuscript is available as a preprint at
\href{https://doi.org/10.5281/zenodo.19174380}{doi.org/10.5281/zenodo.19174380}
(concept DOI: \href{https://doi.org/10.5281/zenodo.19174379}{doi.org/10.5281/zenodo.19174379}).
Companion papers: RAGC~\href{https://doi.org/10.5281/zenodo.19197999}{doi.org/10.5281/zenodo.19197999}, EB~\href{https://doi.org/10.5281/zenodo.19198120}{doi.org/10.5281/zenodo.19198120}, RAHC~\href{https://doi.org/10.5281/zenodo.19198171}{doi.org/10.5281/zenodo.19198171}, RACI~\href{https://doi.org/10.5281/zenodo.19198224}{doi.org/10.5281/zenodo.19198224}, RAQI~\href{https://doi.org/10.5281/zenodo.19198285}{doi.org/10.5281/zenodo.19198285}, RADM~\href{https://doi.org/10.5281/zenodo.19198513}{doi.org/10.5281/zenodo.19198513}.

\textbf{Code availability:} The Lean~4 formal proof files
(\texttt{RA\_AQFT\_Proofs\_v10.lean} and \texttt{RA\_Proofs\_Lean4.lean}),
Python computational scripts, and all companion paper source files are publicly
available at \href{https://github.com/jsandeman/Relational-Actualism}{github.com/jsandeman/Relational-Actualism}.

\textbf{Data availability:} This paper presents theoretical work. No datasets
were generated or analysed. All formal verification results are contained in the
Lean~4 source files available at the repository above.

% ─────────────────────────────────────────────────────────────────────────────
\bibliographystyle{plainnat}
\begin{thebibliography}{99}

\bibitem[Sandeman(2026RACL)]{Sandeman2026RACL}
J.~F. Sandeman.
\newblock {Relational Actualism: The Einstein Field Equations as the Unique
  Macroscopic Limit of the Actualization Graph (RACL)}.
\newblock \textit{Preprint}, 2026.
\newblock \href{https://doi.org/10.5281/zenodo.19270020}{doi.org/10.5281/zenodo.19270020}.


\bibitem[Sandeman(2026RAEB)]{Sandeman2026RAEB}
J.~F. Sandeman.
\newblock {Relational Actualism: The Engine of Becoming (RAEB)}.
\newblock \textit{Preprint}, 2026.
\newblock \href{https://doi.org/10.5281/zenodo.19198120}{doi.org/10.5281/zenodo.19198120}.


\bibitem[Sandeman(2026RATM)]{Sandeman2026RATM}
J.~F. Sandeman.
\newblock {Relational Actualism: The Topology of Matter (RATM)}.
\newblock \textit{Preprint}, 2026.
\newblock \href{https://doi.org/10.5281/zenodo.19265651}{doi.org/10.5281/zenodo.19265651}.


\bibitem[Rovelli(1996)]{Rovelli1996}
Rovelli, C. (1996).
\newblock Relational quantum mechanics.
\newblock \emph{International Journal of Theoretical Physics}, 35(8):1637--1678.

\bibitem[Hartle and Hawking(1983)]{HartleHawking1983}
Hartle, J.~B. and Hawking, S.~W. (1983).
\newblock Wave function of the universe.
\newblock \emph{Physical Review D}, 28:2960--2975.

\bibitem[Zurek(1981)]{Zurek1981}
Zurek, W.~H. (1981).
\newblock Pointer basis of quantum apparatus: into what mixture does the wave packet collapse?
\newblock \emph{Physical Review D}, 24(6):1516--1525.

\bibitem[Zurek(2003)]{Zurek2003}
Zurek, W.~H. (2003).
\newblock Decoherence, einselection, and the quantum origins of the classical.
\newblock \emph{Reviews of Modern Physics}, 75(3):715--775.

\bibitem[Joos and Zeh(1985)]{JoosZeh1985}
Joos, E. and Zeh, H.~D. (1985).
\newblock The emergence of classical properties through interaction with the environment.
\newblock \emph{Zeitschrift f\"{u}r Physik B}, 59(2):223--243.

\bibitem[Minev et~al.(2019)]{Minev2019}
Minev, Z.~K., Mundhada, S.~O., Shankar, S., Reinhold, P.,
  Guti\'{e}rrez-J\'{a}uregui, R., Schoelkopf, R.~J., Mirrahimi, M.,
  Carmichael, H.~J., and Devoret, M.~H. (2019).
\newblock To catch and reverse a quantum jump mid-flight.
\newblock \emph{Nature}, 570:200--204.

\bibitem[Cohen et~al.(2019)]{Cohen2019}
Cohen, E., Larocque, H., Bouchard, F., Nejadsattari, F., Gefen, Y., and
  Karimi, E. (2019).
\newblock Geometric phase from Aharonov--Bohm to Pancharatnam--Berry and beyond.
\newblock \emph{Nature Reviews Physics}, 1(7):437--449.

\bibitem[Hackerm\"{u}ller et~al.(2004)]{Hackermuller2004}
Hackerm\"{u}ller, L., Hornberger, K., Brezger, B., Zeilinger, A., and Arndt, M.
  (2004).
\newblock Decoherence of matter waves by thermal emission of radiation.
\newblock \emph{Nature}, 427:711--714.

\bibitem[Hornberger et~al.(2003)]{Hornberger2003}
Hornberger, K., Uttenthaler, S., Brezger, B., Hackerm\"{u}ller, L., Arndt, M.,
  and Zeilinger, A. (2003).
\newblock Collisional decoherence observed in matter wave interferometry.
\newblock \emph{Physical Review Letters}, 90:160401.

\bibitem[Scully and Dr\"{u}hl(1982)]{ScullyDruhl1982}
Scully, M.~O. and Dr\"{u}hl, K. (1982).
\newblock Quantum eraser: a proposed photon correlation experiment concerning
  observation and `delayed choice' in quantum mechanics.
\newblock \emph{Physical Review A}, 25(4):2208--2213.

\bibitem[Everett(1957)]{Everett1957}
Everett, H. (1957).
\newblock Relative state formulation of quantum mechanics.
\newblock \emph{Reviews of Modern Physics}, 29(3):454--462.

\bibitem[Saunders et~al.(2010)]{SaundersEtAl2010}
Saunders, S., Barrett, J., Kent, A., and Wallace, D., editors (2010).
\newblock \emph{Many Worlds? Everett, Quantum Theory, and Reality}.
\newblock Oxford University Press.

\bibitem[Lieb and Robinson(1972)]{LiebRobinson1972}
Lieb, E.~H. and Robinson, D.~W. (1972).
\newblock The finite group velocity of quantum spin systems.
\newblock \emph{Communications in Mathematical Physics}, 28(3):251--257.

\bibitem[Pazourek et~al.(2015)]{Pazourek2015}
Pazourek, R., Nagele, S., and Burgd\"{o}rfer, J. (2015).
\newblock Attosecond chronoscopy of photoemission.
\newblock \emph{Reviews of Modern Physics}, 87:765--802.

\bibitem[Bassi et~al.(2013)]{Bassi2013}
Bassi, A., Lochan, K., Satin, S., Singh, T.~P., and Ulbricht, H. (2013).
\newblock Models of wave-function collapse, underlying theories, and experimental
  tests.
\newblock \emph{Reviews of Modern Physics}, 85(2):471--527.

\bibitem[Haag and Kastler(1964)]{HaagKastler1964}
Haag, R. and Kastler, D. (1964).
\newblock An algebraic approach to quantum field theory.
\newblock \emph{Journal of Mathematical Physics}, 5(7):848--861.

\bibitem[Sandeman(2026d)]{Sandeman2026RAHC}
J.~F.~Sandeman.
\newblock Algorithmic Causal Dynamics: A Unified Framework for Emergence and
  Complexity in Relational Actualism ({RAHC}).
\newblock Preprint, 2026.
\newblock \href{https://doi.org/10.5281/zenodo.19198171}{doi.org/10.5281/zenodo.19198171}.

\bibitem[Sandeman(2026b)]{Sandeman2026RAGC}
J.~F.~Sandeman.
\newblock Relational Actualism: Gravity and Cosmology ({RAGC}).
\newblock Preprint, 2026.
\newblock \href{https://doi.org/10.5281/zenodo.19197999}{doi.org/10.5281/zenodo.19197999}.

\bibitem[Sandeman(2026c)]{Sandeman2026RACI}
J.~F.~Sandeman.
\newblock Relational Actualism: Complexity and Intelligence ({RACI}).
\newblock Preprint, 2026.
\newblock \href{https://doi.org/10.5281/zenodo.19198224}{doi.org/10.5281/zenodo.19198224}.


\bibitem[Meiburg et~al.(2025)]{Meiburg2025LQI}
A.~Meiburg, L.~A.~Lessa, and R.~R.~Soldati.
\newblock A formalization of the generalized quantum {Stein's} lemma in {Lean}.
\newblock \textit{arXiv}:2510.08672, 2025.
\newblock Lean-QuantumInfo library: \url{https://github.com/Timeroot/Lean-QuantumInfo}.

\bibitem[Sandeman(2026RASM)]{Sandeman2026RASM}
J.~F.~Sandeman.
\newblock Relational Actualism and the Standard Model ({RASM}).
\newblock Preprint, 2026.
\newblock \href{https://doi.org/10.5281/zenodo.19229335}{doi.org/10.5281/zenodo.19229335}.

\bibitem[Sandeman(2026BDG)]{Sandeman2026BDG}
J.~F.~Sandeman.
\newblock Two Standard Model Coupling Constants from the Unique 4-Dimensional
  BDG Integers.
\newblock Preprint, 2026.
\newblock \href{https://doi.org/10.5281/zenodo.19324790}{doi.org/10.5281/zenodo.19324790}.

\end{thebibliography}

\end{document}
